\section{Results}
\label{sec:results}

\subsection{Exact BNN predicate and discovered evaluator}

Table~\ref{tab:bnn-truth} gives the complete predicate extension supplied to the search.  Only basis index one, written in the human ordering $x_1x_2x_3=100$, is marked because $q_0$ stores $x_1$ as the computational-basis least-significant bit.

\begin{table}[t]
\caption{Exact truth table of the verification predicate.}
\label{tab:bnn-truth}
\centering
\begin{tabular}{ccccccccc}
\toprule
$x_1x_2x_3$ & 000&100&010&110&001&101&011&111\\
\midrule
$g(x)$ &0&1&0&0&0&0&0&0\\
$(-1)^{g(x)}$ &1&$-1$&1&1&1&1&1&1\\
\bottomrule
\end{tabular}
\end{table}

The frozen hierarchy reaches an exact evaluator after 94 outer allocations and 352 attempted macro continuations.  The discovered macro witness is
\begin{equation}
\begin{aligned}
&\operatorname{TOF}_{0,1\to4};\;
\CNOT_{0\to4};\;
\CNOT_{4\to3};\\
&\operatorname{TOF}_{2,4\to3};\;
\CNOT_{0\to4};\;
\operatorname{TOF}_{0,1\to4}.
\end{aligned}
\label{eq:found-macro-witness}
\end{equation}
As derived in Section~\ref{sec:boolean-oracle}, this computes $f=x_1\bar x_2\bar x_3$ and restores $q_4$ exactly.  It is not stored in the target object and was not part of the four-function training curriculum.

\subsection{Native Clifford+$T$ certificate}

Table~\ref{tab:oracle-resources} reports resources after exact native lowering.  The evaluator is lowered once, $Z_f=S_f^2$ is inserted, and the exact inverse evaluator is appended.  Consequently the complete phase oracle approximately doubles the evaluator resources, with two additional single-qubit Clifford gates.

\begin{table}[t]
\caption{Certified native resources for the BNN evaluator and phase oracle.}
\label{tab:oracle-resources}
\centering
\begin{tabular}{lrrrr}
\toprule
Circuit & Native gates & $T/T^\dagger$ & CNOT & Depth\\
\midrule
Reversible evaluator $U_g$ & 48 & 21 & 21 & 38\\
Phase oracle $U_g^\dagger Z_fU_g$ & 98 & 42 & 42 & 76\\
\bottomrule
\end{tabular}
\end{table}

The independent exact-phase certifier obtains
\begin{equation}
\Delta_{\mathrm{exact}}=1.224\times10^{-15},
\qquad
\Delta_{\mathrm{proj}}=1.222\times10^{-15},
\end{equation}
and
\begin{equation}
\ell_A=3.265\times10^{-32}.
\end{equation}
These values are numerical roundoff at the scale of the dense calculation.  The flag and work qubit are restored on all eight data basis states, and linearity extends the result to arbitrary input superpositions.

\subsection{Search controls and computational work}

All three frozen procedures synthesize the same six-macro evaluator and pass the same independent phase-oracle certificate.  Table~\ref{tab:oracle-controls} reports medians over five timing repetitions.

\begin{table*}[t]
\caption{Frozen evaluator-search results for the BNN phase oracle.  Timings are environment-specific medians; all methods use identical exact transitions and certification.}
\label{tab:oracle-controls}
\centering
\begin{tabular}{lrrrrr}
\toprule
Method & Certified & Wall time (s) & Macro edges & Outer allocations & Peak frontier\\
\midrule
Hierarchical outer SARSA + inner LinUCB & yes & 0.10329 & 352 & 94 & 236\\
Deferred fixed continuation order & yes & 0.09534 & 371 & 100 & 252\\
Exact mapping target potential & yes & 0.03847 & 180 & 48 & 139\\
\bottomrule
\end{tabular}
\end{table*}

The inner policy reduces attempted edges by $5.1\%$ and the peak frontier by $6.3\%$ relative to fixed continuation order, but its scoring overhead makes it approximately $8.3\%$ slower on this small instance.  The exact target-potential control is both faster and less expensive here.  The result therefore establishes that the current learned hierarchy can synthesize the oracle, not that it is the best scheduler for a uniquely marked three-bit predicate.

\begin{figure*}[t]
\centering
\begin{tikzpicture}
\begin{groupplot}[
  group style={group size=2 by 1,horizontal sep=1.35cm},
  width=0.45\textwidth,height=0.29\textwidth,
  ybar,
  symbolic x coords={Deferred fixed,Exact potential,SARSA+LinUCB},
  xtick=data,xticklabel style={rotate=18,anchor=east,font=\scriptsize},
  tick label style={font=\scriptsize},label style={font=\footnotesize},
  title style={font=\footnotesize},
  enlarge x limits=0.22,
  grid=major]
\nextgroupplot[ylabel={Median wall time (s)},title={Implementation time},ymin=0]
\addplot table[x=label,y=median_wall_seconds,col sep=comma]{data/oracle_method_summary.csv};
\nextgroupplot[ylabel={Exact macro edges},title={Search work},ymin=0]
\addplot table[x=label,y=macro_edges,col sep=comma]{data/oracle_method_summary.csv};
\end{groupplot}
\end{tikzpicture}
\caption{BNN evaluator-search controls.  The learned inner policy performs slightly less exact work than fixed ordering, but the exact target-potential scheduler is strongest on this small truth-table instance.  The plots are generated by PGFPlots directly from the packaged CSV data.}
\label{fig:oracle-controls}
\end{figure*}

\subsection{Search trajectory}

The exact mapping distance reaches zero at allocation 94.  The frontier grows to 236 active records, illustrating why deferred edge scheduling and exact mapping/Pareto collapse remain useful even for three Boolean inputs.

\begin{figure*}[t]
\centering
\begin{tikzpicture}
\begin{groupplot}[
  group style={group size=2 by 1,horizontal sep=1.35cm},
  width=0.45\textwidth,height=0.28\textwidth,
  xlabel={Outer allocation},
  tick label style={font=\scriptsize},label style={font=\footnotesize},
  title style={font=\footnotesize},grid=major]
\nextgroupplot[ylabel={Best exact mapping distance},title={Target progress}]
\addplot table[x=allocation,y=best_distance,col sep=comma]{data/oracle_search_trace.csv};
\nextgroupplot[ylabel={Active frontier records},title={Frontier growth}]
\addplot table[x=allocation,y=frontier_size,col sep=comma]{data/oracle_search_trace.csv};
\end{groupplot}
\end{tikzpicture}
\caption{Trace of the representative hierarchical run.  Nonmonotone changes are expected because the outer policy may allocate effort to a different frontier record and because compute--uncompute routes need not decrease a single local potential at every step.}
\label{fig:oracle-trace}
\end{figure*}

\subsection{Restricted QFT-3 and unrestricted-search boundary}

The restricted QFT-3 generator produces a certified 35-gate native circuit with 15 $T/T^\dagger$ gates, 13 CNOTs, and depth 26.  Its exact-isometry error is $1.38\times10^{-15}$ and leakage is $3.83\times10^{-33}$.  This validates clean-ancilla controlled-phase construction and exact native lowering.

The generic unrestricted four-physical-qubit search does not reach QFT-3 within its 1.5-second, 585-edge probe.  The two observations are not contradictory: the restricted generator is told the QFT block structure, whereas the generic search must discover a 35-gate mixed circuit from a 32-way native branching grammar.  The paper therefore treats QFT-3 as an integration benchmark and the BNN truth-table oracle as the positive specification-level synthesis result.

\subsection{Supporting canonicalization and mixed-gate validation}

The strengthened canonicalization audit records no false merge in the bounded one- and two-qubit enumerations and approximately halves unresolved duplicate excess.  Separate four--six qubit mixed Clifford+$T$ crossover tests show that deferred hierarchical scheduling can reduce native continuation work when per-node branching is larger.  These results support the semantic and computational components reused by the oracle pipeline, although they do not replace the oracle-specific controls above.
